% SGA 3, Expose VI_B, Loop 1 English reconstruction.
% Scope: Proposition 7.4 and its complete proof.
% Authority: Polo--Gille Exp6B-13oct24.pdf, component-local pages
% 54--55 (combined-reader pages 395--396; re-edition headers
% 384--385; original-margin printed pages 367--369).
% The preceding tail of 7.3 on local page 54 is excluded.
% Lemma 7.5 begins on local page 56 and is excluded.
% OCR and Reinhold English Markdown are comparison witnesses only.
% This unit contains no diagrams.

\medskip
\phantomsection\label{prop:sga3-VIB-7.4}%
\noindent\textbf{命题 7.4.---}
\emph{设 \(G\) 为局部有限型 \(k\)-群,
\(X\) 为有限型 \(k\)-概形, 且几何
约化并几何连通, \(f:X\to G\) 为
\(k\)-态射, 使得 \(f(X)\) 包含 \(G\) 的单位元
\(e\). 则采用
\hyperref[concl:sga3-VIB-7.1-ii]{7.1\textup{(ii)}} 的记号,
存在一个整数 \(N\), 使得在集合意义下,}
\phantomsection\label{form:sga3-VIB-7.4-stabilization}%
\[
  f^N(X^N)=\Gamma_G(f).
\]

\emph{证明.}
由 \hyperref[concl:sga3-VIB-7.1-iii]{7.1\textup{(iii)}} 和
EGA~IV\(_2\), 2.6.1, 可以假设 \(k\) 代数
闭. 由推论~\hyperref[cor:sga3-VIB-7.2.1]{7.2.1}, 可以假设
\(G=G^0\). 最后, 采用
\hyperref[rem:sga3-VIB-7.2-v]{7.2\textup{(v)}} 的记号, 只需
证明存在一个整数 \(N\), 使得
\[
  f'^{\,N}(X'^{\,N})=\Gamma_G(f).
\]

\emph{第一种情形.}
假设 \(X\) 不可约. 则不可约闭
子集
\[
  \overline{f'^{\,n}(X'^{\,n})}
\]
在空间 \(G\) 中构成一个递增序列; 此空间
为 Noether 空间, 因为 \(G=G^0\) 在 \(k\) 上有限型
(\hyperref[prop:sga3-VIA-2.4]{VI\(_A\), 2.4}). 因此此
序列最终稳定, 存在一个整数 \(m\), 使得
\[
  \overline{f'^{\,m}(X'^{\,m})}=\Gamma_G(f).
\]

此外, 因为 \(X\) 和 \(G\) 在 \(k\) 上有限型,
态射 \(f'^{\,n}\) 在 \(k\) 上有限型.
因此, \(f'^{\,m}(X'^{\,m})\) 在
\(G\) 中可构造 (EGA~IV\(_1\), 1.8.5), 从而包含其闭包
\(\Gamma_G(f)\) 的一个开子集 \(U\)
(EGA~0\(_{\mathrm{IV}}\), 9.2.3). 由
\hyperref[sga3:VIA:0.5]{VI\(_A\), 0.5}, 此时有
\phantomsection\label{form:sga3-VIB-7.4-first-case-chain}%
\[
  \Gamma_G(f)
  \subset U\cdot U
  \subset f'^{\,2m}(X'^{\,2m})
  \subset \Gamma_G(f),
\]
从而
\[
  f'^{\,2m}(X'^{\,2m})=\Gamma_G(f).
\]

\emph{第二种情形.}
假设 \(X\) 恰有两个不可约分支
\(A_1\) 和 \(A_2\). 由于 \(X\) 连通且 \(k\)
代数闭, 存在
\[
  a\in(A_1\cap A_2)(k).
\]
因此四个不可约子集
\[
  A_i\times_kA_j
  \qquad (i,j=1,2)
\]
覆盖 \(X\times_kX\), 且其中每个在态射
\[
  f'^{\,1}=\mu\circ\bigl(f\times_k(c\circ f)\bigr)
\]
下的像均包含 \(e\). 用 \(f'^{\,1}_{ij}\) 表示
\(f'^{\,1}\) 在约化子概形 \(A_i\times_kA_j\) 上的限制,
并令
\[
\begin{aligned}
  Y={}&
  (A_1\times_kA_1)
  \times_k(A_1\times_kA_2)
  \times_k(A_2\times_kA_2)
  \times_k(A_2\times_kA_1),\\
  g={}&
  \mu\circ\Bigl(
    \mu\circ
      (f'^{\,1}_{11}\times_k f'^{\,1}_{12})
    \times_k
    \mu\circ
      (f'^{\,1}_{22}\times_k f'^{\,1}_{21})
  \Bigr).
\end{aligned}
\]
概形 \(Y\) 不可约、约化且有限型.
因此由第一种情形, 存在一个整数 \(m\),
使得
\[
  g'^{\,m}(Y'^{\,m})=\Gamma_G(g).
\]
但对每个 \(n\geq1\),\begingroup
\renewcommand{\thefootnote}{\fnsymbol{footnote}}%
\footnote[1]{\label{note:sga3-VIB-7.4-Yprime}译者注:
法文底本将中间项印作 \(g'^{\,n}(Y^{\,n})\).
由于 \(g'^{\,n}\) 的源为 \(Y'^{\,n}\),
此处补回了撇号; 比较上面的等式
\(g'^{\,m}(Y'^{\,m})=\Gamma_G(g)\).}%
\endgroup
\phantomsection\label{form:sga3-VIB-7.4-two-component-chain}%
\[
  f'^{\,n}(X'^{\,n})
  \subset g'^{\,n}(Y'^{\,n})
  \subset f'^{\,4n}(X'^{\,4n}).
\]
由此得到
\[
  \Gamma_G(f)=\Gamma_G(g)
  \qquad\text{且}\qquad
  f'^{\,4m}(X'^{\,4m})=\Gamma_G(f).
\]

\emph{一般情形.}
对 \(X\) 的不可约分支数作归纳;
此数有限, 因为 \(X\) 在 \(k\) 上有限
型, 从而为 Noether 概形. 假设当 \(X\) 至多有 \(r-1\) 个不可约
分支时, 命题已证, 并假设 \(X\) 有 \(r\) 个这样的分支,
记为 \(A_1,\ldots,A_r\). 单位元 \(e\) 属于
某个 \(A_i\) 的像; 例如假设
\[
  e\in f(A_1).
\]
令
\[
  X_r=A_1\cup\cdots\cup A_{r-1},
  \qquad
  Y=\Gamma_G(f_r),
\]
其中 \(f_r\) 表示 \(f\) 在约化
子概形 \(X_r\) 上的限制. 选取 \(A_i\) 的编号,
使 \(X_r\) 连通; 这总是可能的. 则由
推论~\hyperref[cor:sga3-VIB-7.2.1]{7.2.1}, \(Y\)
为 \(G\) 的闭约化不可约子群.

令
\[
  Z=\overline{f(A_r)}
\]
并设 \(T=Y\cup Z\), 赋予其闭约化
子概形结构; 设 \(i:T\hookrightarrow G\) 为
包含态射. 由
\hyperref[rem:sga3-VIB-7.2-ii]{7.2\textup{(ii)}} 立即得到
\[
  \Gamma_G(i)=\Gamma_G(f).
\]
此外, \(T\) 连通: 由于 \(X\) 连通,
\(X_r\) 和 \(A_r\) 有一个公共点 \(a\), 而 \(Y\) 和
\(Z\) 有公共点 \(f(a)\). 另外, \(e\in T\), 且
\(T\) 至多有两个不可约分支, 因为 \(Y\)
和 \(Z\) 都不可约.

由归纳假设, 存在一个整数 \(m\),
使得
\[
  f_r'^{\,m}(X_r'^{\,m})
  =\Gamma_G(f_r)=Y.
\]
由于 \(X=X_r\cup A_r\), 且因
\(e\in f(A_1)\), 有
\[
  Z=\overline{f(A_r)}
  \subset\overline{f'^{\,m}(X'^{\,m})}
\]
故得到
\phantomsection\label{form:sga3-VIB-7.4-general-case-inclusion}%
\[
  f(X)
  \subset Y\cup Z
  \subset\overline{f'^{\,m}(X'^{\,m})}.
\]

另一方面, 因为 \(T\) 至多有两个不可约
分支, 第二种情形给出一个整数 \(p\), 使得
\[
  i'^{\,p}(T'^{\,p})
  =\Gamma_G(i)=\Gamma_G(f).
\]
又有
\[
  T\subset\overline{f'^{\,m}(X'^{\,m})},
\]
从而
\[
  \overline{f'^{\,mp}(X'^{\,mp})}=\Gamma_G(f).
\]
与第一种情形中一样, 此最后等式蕴涵
\[
  f'^{\,2mp}(X'^{\,2mp})=\Gamma_G(f).
\]
这就完成了证明.

% 翻译统计: 源行数 213; 输出行数 212 (含本统计行); 行内/行间数学片段 109; 翻译占位 0. 底本疑点另报主线程.

% 独立复核检查点：已逐段对照英文源文，覆盖正文、假设、证明、脚注与源注。
% 行内/陈列数学及全部引用键机械对照通过；仅统一群的invariant为“正规”，保留源文疑点另报。
